\section{Introduction}
\label{sec:intro}

% P1: task and why token cost matters
Language models do not read characters or words. A tokenizer first splits the text into
subword units, and the cost of a request, its latency and the amount of text that fits in
the context window are all counted in tokens. Byte-pair encoding (BPE)
\citep{gage1994,sennrich2016}, the most widely used tokenizer, builds its vocabulary by
repeatedly merging the most frequent pair of adjacent symbols in a training corpus. Since
the merges follow frequency, they follow the languages that dominate the corpus. The same
content then needs many more tokens in some languages than in others, with differences of
up to 15 times for the tokenizers of deployed models \citep{petrov2023}. Consider a Hindi
speaker and an English speaker sending the same request to a model whose tokenizer was
trained on an English-heavy mixture. In our setting, the Hindi request costs 21\% more
tokens (35.0 vs.\ 29.0 per sentence), so its user pays more, waits longer and fits less
text into the same context window.

% P2: existing methods sit at two ends, the space between is unprincipled
Existing tokenizers answer this disparity at two opposite ends. Classical BPE minimises the
total number of tokens and accepts whatever disparity results. Parity-aware BPE
\citep{foroutan2026} gives every merge to the language that is currently worst off, which
equalises the cost across languages at some loss of total compression. To reach points in
between, \citet{foroutan2026} propose a hybrid schedule that switches from the classical
rule to the parity rule after a fixed share of merges, and \citet{kanjirangat2026} extend
the parity rule to the $N$ worst-off languages. However, both are heuristic interpolations:
neither is derived from an objective that describes the trade-off, so a practitioner cannot
state what the intermediate tokenizers optimise or choose a point on the trade-off
deliberately. Moreover, Parity-aware BPE controls only the worst-off language, and
\citet{lee2026} show that it becomes less equitable as the vocabulary grows, since nothing
stops the dominant language from acquiring more merges.

% P3: our proposal
To address this, we treat vocabulary learning as a multi-objective problem over the vector
of per-language token costs. We scalarise this vector with a weighted power mean of
exponent $p$, and show that the greedy step of this objective is a pair count in which each
language is weighted by its current cost raised to the power $p-1$
(Section~\ref{sec:method}). At $p=1$ the rule is exactly classical BPE, and as
$p\to\infty$ it becomes the rule of Parity-aware BPE. We call the tokenizers in between
\emph{$p$-fair BPE}. Since every language keeps a non-zero weight for finite $p$, the rule
favours expensive languages without ignoring the rest. Additionally, we prove that the
score of an existing pair never increases during training, which allows exact training with
a lazy priority queue \citep{minoux1978} at about 20 times the speed of full rescoring.

% P4: what we find
We evaluate $p$-fair BPE on Hindi, Marathi, Gujarati and English, with five samples of
training text and a paired bootstrap on FLORES-200 \citep{nllb2022}. Raising $p$ moves the
tokenizer monotonically from the classical end to the parity end in five of five samples.
At $p=8$, the Gini coefficient of per-language cost falls by 40\% for a 0.6\% loss in
compression of the training text. The resulting front coincides with that of the hybrid
schedule, so $p$-fair BPE matches it and does not dominate it. Nevertheless, the two methods
distribute the cost differently: for every value of $p$ we test, no Indian language costs
more than under classical BPE, whereas the hybrid raises the cost of Marathi (33.26 vs.\
32.47 tokens per sentence at 75\% classical merges). Our experiments also expose two factors
that decide the outcome more than the merge rule does. First, measuring cost per byte
instead of per parallel sentence inverts the parity rule across scripts, sending 99\% of
its merges to English. Second, the disparity depends strongly on data balance, as the Gini
coefficient of classical BPE ranges from 0.019 with equal data to 0.051 with eight times
as much English.

\noindent\textbf{Contributions.} Our main contributions are:
\begin{itemize}
  \item We formulate fair vocabulary learning as the minimisation of a power mean of
  per-language token costs, and show that its greedy step contains classical BPE ($p=1$)
  and Parity-aware BPE ($p\to\infty$) as limits of one rule.
  \item We prove that pair scores are non-increasing under this rule, which gives an exact
  lazy-evaluation trainer about 20 times faster than full rescoring.
  \item We trace the compression--parity front on Hindi, Marathi, Gujarati and English with
  error bars over five training samples, and show that $p$-fair BPE matches the hybrid
  front while never raising the cost of an under-served language above its classical cost.
  \item We show that the cost unit and the training data balance decide the size and even
  the direction of the disparity, and give practical guidance for both.
\end{itemize}